\section{Experimental Setup}
\label{sec:experiments}

\subsection{Evaluation Metrics}
\textbf{T1, fault type.} We report accuracy, macro-averaged $F_1$ (the unweighted mean over the 7 classes, so rare and hard classes count as much as easy ones) and $7\times7$ confusion matrices.

\textbf{T2, fault location.} The metrics are computed over every trace whose true class is not \texttt{none}, whatever class was predicted. A missed fault therefore still counts, with whatever location the method outputs; the conventional detector outputs $d/l = 0.5$ when it finds no edge. We report the mean absolute error $\mathrm{MAE} = \frac{1}{N}\sum |\hat d_i - d_i|$ in metres, the median error, the 90th-percentile error $P_{90}$ (which captures the worst cases), and the share of faults located within 1~m.

\textbf{T3, load.} These metrics are computed only where the load is observable, i.e.\ the fault is not a short or open. We report accuracy and macro-$F_1$ over the 6 load classes, and the mean absolute percentage error (MAPE) of $R_L$, of $L = \tau (R_L + Z_0)$ for true $RL$ loads, and of $C = \tau (R_L + Z_0)/(R_L Z_0)$ for true $RC$ loads.

\subsection{Experiments}
Table~\ref{tab:experiments} lists the eight experiments. The test split is used only for E1--E4, E7 and E8. Every choice of threshold, model or stopping epoch was made on the validation split.

\begin{table}[!htbp]
\centering
\caption{The eight experiments. ``Seed 0'' means the models trained with seed 0; ``3 seeds'' gives mean $\pm$ standard deviation over seeds 0--2.}
\label{tab:experiments}
\small
\renewcommand{\arraystretch}{1.2}
\begin{tabular}{@{}l p{0.60\textwidth} l@{}}
\toprule
\textbf{ID} & \textbf{Question and set-up} & \textbf{Data / models} \\
\midrule
E1 & Main comparison of the three methods at 30~dB, with one fixed noise realisation. & Test, 3 seeds \\
E2 & Noise robustness at 60, 50, 40, 30, 25, 20 and 15~dB (15 and 60~dB are outside the 20--50~dB training range). & Test, seed 0 \\
E3 & Generalisation to the five OOD subsets of Section~\ref{sec:simulation}. & OOD, seed 0 \\
E4 & Where the errors come from: location error by position and by $|\Gamma|$, and detection recall per $|\Gamma|$ bin. & Test, seed 0 \\
E5 & Data efficiency: training on 5, 10, 25, 50 and 100\% of the training split. & Validation, seed 0 \\
E6 & Ablations: input channels $(\tilde v, d\tilde v/dt, t/T)$ vs.\ $(\tilde v, d\tilde v/dt)$ vs.\ $\tilde v$ alone, and one multi-task network vs.\ three single-task networks. & Validation, seed 0 \\
E7 & What the LightGBM fault classifier relies on: mean $|$SHAP$|$ \cite{lundberg2017unified} on 2,000 test traces. & Test, seed 0 \\
E8 & Physics check on real hardware: velocity factor, damage location and record consistency from echo timing on the WISE-CASING records \cite{nrel_wise_casing}. & Lab data, no ML \\
\bottomrule
\end{tabular}
\end{table}
